% Part: counterfactuals
% Chapter: introduction
% Section: counterfactuals

\documentclass[../../../include/open-logic-section]{subfiles}

\begin{document}

\olfileid{cnt}{int}{cnt}

\olsection{Kontrafaktual}

Wujud konstruksi ``if \dots then \dots'' kang lumrah lan wigati
ing basa Inggris nggunakake wujud subjunktif wektu kepungkur saka
\emph{to be}: ``if it were the case that \dots then it would be the
case that \dots'' (``saupama ... mesthi ...''). Amarga lumrahe
anteseden kondisional kaya mangkono luput, yaiku lelawanan karo
fakta, kondisional iku diarani \emph{kondisional kontrafaktual}
(lan amarga nggunakake wujud subjunktif saka \emph{to be}, uga
diarani \emph{kondisional subjunktif}). Kondisional iku dibedakake
saka kondisional \emph{indikatif}, kang awujud ``if it is the case
that \dots then it is the case that \dots''. Kondisional
kontrafaktual lan indikatif beda syarat bebener. Gatekna tuladha
Adams kang misuwur:
\begin{quote}
  Yen Oswald ora mateni Kennedy, wong liya kang mateni.

  Saupama Oswald ora mateni Kennedy, wong liya mesthi wis mateni.
\end{quote}
Ukara kapisan indikatif, dene kapindho kontrafaktual. Ukara kapisan
cetha bener: kita ngerti yen Presiden John F. Kennedy dipateni
dening \emph{sawijining wong}, lan yen wong iku dudu Lee Harvey
Oswald (kang lelawanan karo Laporan Warren), mula wong liya kang
mateni Kennedy. Ukara kapindho ngandharake prakara kang beda.
Pratelane yaiku saupama Oswald ora mateni Kennedy, yaiku saupama
penembakan ing Dallas bisa dicegah utawa ora kasil, sejarah
sabanjure mesthi lumaku kanthi cara kang ndadekake upaya
pembunuhan liyane kasil. Supaya ukara iku bener, kudu ana
kakuwatan gedhe kang wis komplotan kanggo njamin patine JFK
(kaya kang dipercaya akeh panyengkuyung teori komplotan JFK).

Apa kondisional \emph{indikatif} bisa diwakili kanthi bener dening
kondisional material isih dadi rembugan. Mligine, apa paradoks
kondisional material bisa ``diterangake'' kanthi cara kang cocog
karo kondisional iku menehi syarat bebener kanggo kondisional
indikatif basa Inggris. Kosok balene, ora dadi pasulayan yen
kondisional kontrafaktual ora bisa disimbolake kanthi bener
nganggo kondisional material. Iki cetha amarga, sanajan lumrahe
anteseden kontrafaktual luput, ora kabeh kontrafaktual kanthi
anteseden kang luput iku bener---upamane, yen kita percaya marang
Laporan Warren lan ora ana komplotan kanggo mateni JFK,
kondisional kontrafaktual Adams iku sawijining tuladha.

Kondisional kontrafaktual wigati ing penalaran kausal: sawijining
guna utamane yaiku ngandharake relasi sebab-akibat. Upamane,
nggesek penthol korek njalari penthol iku murub, lan iki bisa
diandharake kanthi ``saupama penthol korek iki digesek, mesthi
bakal murub.'' Kondisional material, lan lumrahe kondisional
indikatif, ora bisa digunakake kanggo ngandharake iki:
``penthol korek digesek $\lif$ penthol korek murub'' bener yen
penthol korek iku ora tau digesek, tanpa nimbang apa kang bakal
kelakon saupama digesek. Malah, ``penthol korek digesek $\lif$
penthol korek malih dadi buntelan kembang'' uga bener yen ora
tau digesek, nanging penthol korek mesthi ora bakal malih dadi
buntelan kembang saupama digesek.

Logika kontrafaktual kang persis bener isih dadi rembugan.
Stalnaker lan Lewis menehi analisis kontrafaktual kang akeh
pengaruhe. Miturut dheweke, kontrafaktual ``saupama~$S$,
mesthi~$T$'' bener yen lan mung yen $T$ bener ing kaanan
kontrafaktual (``donya kang bisa ana'') kang paling cedhak
karo kaanan donya aktual lan ing kono~$S$ bener. Iki diarani
analisis ``ontik'', amarga ngrujuk marang ontologi donya-donya
kang bisa ana. Analisis liyane nggunakake probabilitas
kondisional utawa teori revisi kapercayan. Ana akeh logika
kontrafaktual kang diajokake. Malah ora ana siji logika
kontrafaktual Lewis--Stalnaker: sanajan Stalnaker lan Lewis
menehi panjelasan kanthi arah kang padha, kang ngrujuk marang
donya-donya kang bisa ana lan paling cedhak, asumsi kang
digunakake ngasilake inferensi sah kang beda-beda.

\end{document}
